\section*{अध्याय \ref{ch:g12:space} --- अंतरिक्ष में सदिश, रेखाएँ और समतल}

\begin{solution}{exo:g12:space:1}
$\vect{AB}(2, 1, -1)$ और $\vect{AC}(1, -1, 1)$ समानुपाती नहीं हैं ($\frac21 \neq \frac{1}{-1}$), इसलिए बिंदु संरेख नहीं हैं। रेखा $(AB)$:
\[
(x, y, z) = (1 + 2t,\; t,\; 2 - t), \qquad t \in \R .
\]
\end{solution}

\begin{solution}{exo:g12:space:2}
$2(x - 1) - (y - 1) + 3(z - 0) = 0$, \emph{अर्थात्}
\[
2x - y + 3z - 1 = 0 .
\]
$B(0,2,1)$ के लिए: $0 - 2 + 3 - 1 = 0$: हाँ, $B$ समतल पर है।
\end{solution}

\begin{solution}{exo:g12:space:3}
$\vect{AB}(1,1,0)$, $\vect{AC}(1,0,1)$:
\[
\cos\widehat{A}
= \frac{\vect{AB}\cdot\vect{AC}}{\norm{\vect{AB}}\,\norm{\vect{AC}}}
= \frac{1}{\sqrt2\,\sqrt2} = \frac12,
\]
इसलिए $\widehat A = 60^\circ$।
\end{solution}

\begin{solution}{exo:g12:space:4}
$\mathcal D$ के बिंदु: $(1 + t,\, 2 - t,\, 2t)$। $\mathcal P$ के \omterm{def:g10:algebra:equation}{समीकरण} में रखने पर:
\[
2(1+t) + (2 - t) - 2t + 1 = 5 - t = 0 \iff t = 5 .
\]
अद्वितीय प्रतिच्छेद बिंदु: $(6, -3, 10)$।
\end{solution}

\begin{solution}{exo:g12:space:5}
\omterm{def:g12:space:normal}{अभिलंब सदिश} $(1, -1, 2)$ और $(2, 1, 1)$ \omterm{def:g12:space:collinear}{संरेखी} नहीं हैं, इसलिए दोनों समतल किसी रेखा पर मिलते
हैं। दोनों \omterm{def:g10:algebra:equation}{समीकरण} जोड़ने पर $3x + 3z = 5$, इसलिए $x = \frac53 - z$; फिर पहला \omterm{def:g10:algebra:equation}{समीकरण} $y = x + 2z - 1 = \frac23 + z$ देता है।
$z = t$ के साथ:
\[
(x, y, z) = \left(\frac53 - t,\; \frac23 + t,\; t\right), \qquad t \in \R,
\]
यह \omterm{def:g11:vect:direction}{दिक् सदिश} $(-1, 1, 1)$ वाली रेखा है। (दोनों \omterm{def:g10:algebra:equation}{समीकरण} ज्यों के त्यों संतुष्ट होते
हैं।)
\end{solution}

\begin{solution}{exo:g12:space:6}
दिशाएँ $\vec u_1(1, 2, -1)$ और $\vec u_2(1, 1, 2)$ \omterm{def:g12:space:collinear}{संरेखी} नहीं हैं, इसलिए रेखाएँ \omterm{def:g12:seq:arith-geom}{समांतर} नहीं हैं। प्रतिच्छेद
के लिए
\[
1 + t = 2 + s, \qquad 2t = 1 + s, \qquad 3 - t = 1 + 2s .
\]
चाहिए। पहले दो $t = 1 + s$ और $2(1+s) = 1 + s$ देते हैं, इसलिए $s = -1$, $t = 0$; पर तब तीसरा $3 = -1$
बनता है, जो असंभव है। न कोई उभयनिष्ठ बिंदु, न समांतरता: रेखाएँ विषमतलीय हैं।
\end{solution}

\begin{solution}{exo:g12:space:7}
\emph{1.} $\operatorname{dist} = \dfrac{\abs{2\cdot3 - 1 + 2\cdot1 - 3}}
{\sqrt{4 + 1 + 4}} = \dfrac{4}{3}$।

\emph{2.} $\vec n(2, -1, 2)$ के साथ $H = M_0 + t\,\vec n$, और $t$ ऐसा चुनकर कि $H \in \mathcal P$ हो:
\[
2(3 + 2t) - (1 - t) + 2(1 + 2t) - 3 = 4 + 9t = 0 \iff t = -\frac49 .
\]
अतः $H\left(\frac{19}{9}, \frac{13}{9}, \frac19\right)$, और
\[
M_0H = \norm{t\,\vec n} = \frac49 \times 3 = \frac43,
\]
जो दूरी-सूत्र से मेल खाता है।
\end{solution}

\begin{solution}{exo:g12:space:8}
\omterm{def:g10:coordgeom:system}{निर्देशांक}: $G(1,1,1)$, और समतल $(BDE)$ $B(1,0,0)$, $D(0,1,0)$, $E(0,0,1)$ से होकर जाता है।

\emph{1.} समतल $(BDE)$ का \omterm{def:g10:algebra:equation}{समीकरण} $x + y + z = 1$ है (तीनों बिंदु उसे संतुष्ट करते हैं और वे
संरेख नहीं हैं), इसलिए $\vec n(1,1,1)$ उसका अभिलंब है। चूँकि $\vect{AG} = (1,1,1) = \vec n$ है, इसलिए विकर्ण $(AG)$
समतल $(BDE)$ के \omterm{def:g11:scal:orthogonal}{लांबिक} है।

\emph{2.} $A(0,0,0)$ से दूरी: $\frac{\abs{0 + 0 + 0 - 1}}{\sqrt3} = \frac{1}{\sqrt3} = \frac{\sqrt3}{3}$। रेखा $(AG)$ $(t, t, t)$ है; वह समतल से तब मिलती है जब
$3t = 1$: बिंदु $\left(\frac13, \frac13, \frac13\right)$ पर, जो त्रिभुज $BDE$ का केन्द्रक है।
\end{solution}

\begin{solution}{exo:g12:space:9}
\emph{1.} $\vect{AB}(1, 0, -1)$, $\vect{AC}(-1, 1, 0)$, $\vect{AD}(1, 1, 1)$। यदि $\vect{AD} = a\vect{AB} + b\vect{AC}$ हो, तो \omterm{def:g10:coordgeom:system}{निर्देशांक} $a - b = 1$, $b = 1$, $-a = 1$ देते
हैं: पहले दो $a = 2$ अनिवार्य कर देते हैं, जो $a = -1$ का खंडन है। \omterm{def:g12:space:collinear}{समतलीय} नहीं।

\emph{2.} $\vect{BC}(-2, 1, 1)$, $\vect{BD}(0, 1, 2)$। $\vec n \cdot \vect{BC} = \vec n \cdot \vect{BD} = 0$ हल कीजिए: $b + 2c = 0$ से $b = -2c$; फिर $-2a - 2c + c = 0$ $c = -2a$ देता है।
$a = 1$ के साथ: $\vec n(1, 4, -2)$। $B(2,1,0)$ से होकर समतल:
\[
x + 4y - 2z - 6 = 0 .
\]

\emph{3.} $A(1,1,1)$ से दूरी: $\dfrac{\abs{1 + 4 - 2 - 6}}{\sqrt{21}} = \dfrac{3}{\sqrt{21}}$। $BCD$ का क्षेत्रफल: $\norm{\vect{BC}}^2 = 6$, $\norm{\vect{BD}}^2 = 5$, $\vect{BC}\cdot\vect{BD} = 3$, इसलिए
\[
\text{क्षेत्रफल} = \frac12\sqrt{6 \times 5 - 9} = \frac{\sqrt{21}}{2}.
\]

\emph{4.}
\[
V = \frac13 \times \frac{\sqrt{21}}{2} \times \frac{3}{\sqrt{21}}
= \frac12 .
\]
\end{solution}

\begin{solution}{pb:g12:space:1}
\textbf{1.} समतल: $x + y + z = 1$। रेखा: $(t, t, t)$। प्रतिच्छेद: $3t = 1$: $\left(\frac13, \frac13, \frac13\right)$।

\textbf{2.} $\dfrac{\abs{0 + 0 + 0 - 1}}{\sqrt3} =
\dfrac{1}{\sqrt3} = \dfrac{\sqrt3}{3}$।

\textbf{3.} $1 - 1 + 0 = 0$: \omterm{def:g11:scal:orthogonal}{लांबिक}।

\textbf{4.} वही अभिलंब $(1,1,1)$, पर $\frac52 \neq 1$: कड़ाई से \omterm{def:g12:seq:arith-geom}{समांतर} समतल। रेखा: दिशा
$(1,1,0) \cdot (1,1,1) = 2 \neq 0$: समतल के \omterm{def:g12:seq:arith-geom}{समांतर} नहीं, इसलिए वह उसे एक बिंदु पर भेदती है।

\textbf{5.} \omterm{prop:g10:coordgeom:midpoint}{मध्यबिंदु}: $\left(1, \frac12, 0\right)$, $\left(\frac12, 1, 0\right)$, $\left(0, 1, \frac12\right)$, $\left(0, \frac12, 1\right)$, $\left(\frac12, 0, 1\right)$, $\left(1, 0, \frac12\right)$: हर एक के
\omterm{def:g10:coordgeom:system}{निर्देशांकों} का योग $\frac32$ है।

\textbf{6.} $P$ का अभिलंब $(1,1,1)$ है, जो $(AG)$ की दिशा है: लंब। केंद्र $\left(\frac12,\frac12,\frac12\right)$
के \omterm{def:g10:coordgeom:system}{निर्देशांकों} का योग $\frac32$ है: यानी वह $P$ पर है।

\textbf{7.} क्रमागत मध्यबिंदुओं का अंतर $\left(-\frac12, \frac12, 0\right)$ जैसे \omterm{def:g11:vect:vector}{सदिश} हैं: हर बार लंबाई
$\frac{\sqrt2}{2}$ --- छह बराबर भुजाएँ।

\textbf{8.} $\left(\frac12, 1, 0\right)$ पर: कोर-सदिश $\left(\frac12, -\frac12, 0\right)$ और $\left(-\frac12, 0, \frac12\right)$: बिंदु-गुणनफल $-\frac14$, परिमाण
$\frac{\sqrt2}{2}$: $\cos = -\frac12$: कोण $120^\circ$। बराबर भुजाएँ, बराबर $120^\circ$ कोण: एक सम षट्भुज।

\textbf{9.} परिमाप $6 \times \frac{\sqrt2}{2} = 3\sqrt2
\approx 4.24$। क्षेत्रफल $6 \times \frac{\sqrt3}{4} \times
\frac12 = \frac{3\sqrt3}{4} \approx 1.30$: किसी फलक (क्षेत्रफल $1$) से बड़ा
--- यानी डिब्बे की किसी भी भुजा से बड़ी काट; केवल विकर्ण आयत ($\sqrt2 \approx 1.41$) उसे हरा
पाता है।

\textbf{10.} $(t,t,t)$ $P$ से $t = \frac12$ पर मिलता है: यानी केंद्र --- जो $[AG]$ का
\omterm{prop:g10:coordgeom:midpoint}{मध्यबिंदु} है ($A$ $t = 0$ पर, $G$ $t = 1$ पर)।

\textbf{11.} $c = \frac12$ के लिए: समतल $A$ के कोने पर तीनों कोरों को $\left(\frac12,0,0\right)$, $\left(0,\frac12,0\right)$,
$\left(0,0,\frac12\right)$ पर काटता है: भुजा $\frac{\sqrt2}{2}$ वाला समबाहु त्रिभुज। $c$ बढ़ने पर त्रिभुज
बड़ा होता जाता है, उसके कोने कटकर षट्भुज बन जाते हैं (जो ठीक $c = \frac32$ पर सम है),
और फिर चित्र सममित रूप से सिकुड़कर $G$ के कोने पर त्रिभुज बन जाता है: यही
रत्न-कटाई जैसा रूपांतर है।

\textbf{12.} काट की हर भुजा कटान-समतल का घन के किसी एक \emph{फलक} के साथ
प्रतिच्छेद है, और कोई समतल किसी फलक (सपाट बहुभुज) से अधिक से अधिक एक रेखाखंड
में मिलता है: अतः अधिक से अधिक $6$ भुजाएँ। सात असंभव है; तीन से छह तक सब
होते हैं (प्रश्न 8 और 11 उनमें से दो दिखा देते हैं)।

\textbf{13.} आधार $ABD$: क्षेत्रफल $\frac12$ का समकोण त्रिभुज; ऊँचाई $AE = 1$: आयतन
$\frac13 \times \frac12 \times 1 = \frac16$।

\textbf{14.} $B, D, E, G$ में से हर युग्म ठीक दो \omterm{def:g10:coordgeom:system}{निर्देशांकों} में $1$ से भिन्न है:
छहों दूरियाँ $\sqrt2$ के बराबर --- यानी सम। घन $BDEG$ और $ABDE$ के सर्वांगसम चार
कोने वाले चतुष्फलकों में बँट जाता है: आयतन $1 - 4 \times \frac16 = \frac13$।

\textbf{15.} समतल $(BDE)$: $x + y + z = 1$; केंद्र $\left(\frac12,\frac12,\frac12\right)$: दूरी $\frac{\abs{\frac32 - 1}}{\sqrt3} = \frac{1}{2\sqrt3} =
\frac{\sqrt3}{6}$।

\textbf{16.} $\vect{AG} = (1,1,1)$, $\vect{BH} = (-1,1,1)$: $\cos\theta = \frac{-1 + 1 + 1}{3} =
\frac13$: $\theta \approx 70.5^\circ$, संपूरक $\approx 109.5^\circ$। चार इलेक्ट्रॉन-युग्म
एक-दूसरे को धकेलते हैं और कार्बन के चारों ओर जितना दूर हो सके उतना फैल जाते
हैं: इष्टतम व्यवस्था सम चतुष्फलक के कोने हैं, जिनकी केंद्र-से-शीर्ष दिशाएँ
ठीक इसी कोण पर मिलती हैं --- हीरा प्रश्न 14 का ही रवा है।

\textbf{17.} $P$ पर $(1, t, t)$: $1 + 2t = \frac32$: $t = \frac14$: बिंदु $\left(1, \frac14, \frac14\right)$।

\textbf{18.} $(AB)$: बिंदु $(t, 0, 0)$; $(EG)$: बिंदु $(s, s, 1)$। मिलने के लिए तीसरे
\omterm{def:g10:coordgeom:system}{निर्देशांक} में $0 = 1$ चाहिए: कभी नहीं। \omterm{def:g12:seq:arith-geom}{समांतर} होने के लिए $(1,0,0)$ और $(1,1,0)$ का
\omterm{def:g12:space:collinear}{संरेखी} होना चाहिए: नहीं। न मिलती हैं न \omterm{def:g12:seq:arith-geom}{समांतर} हैं: विषमतलीय --- अलग-अलग मंज़िलों
के दो गलियारे।

\textbf{19.} $P$ पर $B + t(1,1,1) = (1 + t, t, t)$: $1 + 3t = \frac32$: $t = \frac16$: प्रक्षेप $\left(\frac76, \frac16, \frac16\right)$, जिसके \omterm{def:g10:coordgeom:system}{निर्देशांकों}
का योग $\frac32$ है: यानी वह $P$ पर है, जैसा चाहिए था।

\textbf{20.} प्राचलिक रेखाएँ भेदती हैं (प्रश्न 1, 10, 17); \omterm{def:g12:space:normal}{अभिलंब सदिश} कोण
और दूरियाँ नापते हैं (प्रश्न 2, 8, 15, 16); और समतल-समीकरण काटें तराशते हैं
(त्रिभुज-से-षट्भुज वाला रूपांतर)। और खेल के मैदान ने आतिथ्य का पूरा मोल चुका
दिया: वर्गों के डिब्बे में एक सम षट्भुज, कोनों में एक सम चतुष्फलक, हीरे का
कोण, और दो ऐसी रेखाएँ जो एक-दूसरे को अनदेखा कर देती हैं --- वह ज्यामिति जो
केवल अंतरिक्ष ही दे सकता है।
\end{solution}
